\documentclass[9pt,a4paper,twocolumn]{ctexart}

% --- Packages ---
\usepackage[top=1.2cm, bottom=1.2cm, left=1.2cm, right=1.2cm, columnsep=0.8cm]{geometry}
\usepackage{hyperref}
\usepackage{graphicx}
\usepackage{float}
\usepackage{longtable}
\usepackage{tabularx}
\usepackage{booktabs}
\usepackage{enumitem}
\usepackage{xcolor}
\usepackage{fancyhdr}
\usepackage{listings}
\usepackage{fontspec}
\usepackage{titlesec}
\usepackage{caption}

% --- Code block styling ---
\definecolor{codebg}{RGB}{245,245,245}
\definecolor{codeframe}{RGB}{220,220,220}
\definecolor{codekw}{RGB}{0,0,192}
\definecolor{codecomment}{RGB}{0,128,0}
\definecolor{codestring}{RGB}{163,21,21}
\definecolor{codenum}{RGB}{128,0,128}
\definecolor{codepunct}{RGB}{90,90,90}
\definecolor{badgegreen}{RGB}{34,139,34}
\definecolor{badgeblue}{RGB}{30,144,255}
\definecolor{badgeblack}{RGB}{50,50,50}

% --- Difficulty badges (rounded corners via tikz, pre-rendered in saveboxes) ---
\usepackage{tikz}
\newsavebox{\badgechubox}
\savebox{\badgechubox}{\tikz[baseline]{\node[fill=badgegreen,text=white,rounded corners=2.5pt,inner xsep=4pt,inner ysep=1.5pt,font=\sffamily\footnotesize\bfseries]{初级};}}
\newsavebox{\badgezhongbox}
\savebox{\badgezhongbox}{\tikz[baseline]{\node[fill=badgeblue,text=white,rounded corners=2.5pt,inner xsep=4pt,inner ysep=1.5pt,font=\sffamily\footnotesize\bfseries]{中级};}}
\newsavebox{\badgegaobox}
\savebox{\badgegaobox}{\tikz[baseline]{\node[fill=badgeblack,text=white,rounded corners=2.5pt,inner xsep=4pt,inner ysep=1.5pt,font=\sffamily\footnotesize\bfseries]{高级};}}
\newcommand{\lvchu}{\usebox{\badgechubox}}
\newcommand{\lvzhong}{\usebox{\badgezhongbox}}
\newcommand{\lvgao}{\usebox{\badgegaobox}}
\pdfstringdefDisableCommands{%
  \def\lvchu{[初级]}%
  \def\lvzhong{[中级]}%
  \def\lvgao{[高级]}%
}

\lstset{
  basicstyle=\ttfamily\scriptsize,
  keywordstyle=\color{codekw}\bfseries,
  commentstyle=\color{codecomment}\itshape,
  stringstyle=\color{codestring},
  backgroundcolor=\color{codebg},
  frame=single,
  rulecolor=\color{codeframe},
  breaklines=true,
  breakatwhitespace=false,
  breakindent=0pt,
  postbreak=\mbox{\textcolor{red}{$\hookrightarrow$}\space},
  numbers=left,
  numberstyle=\tiny\color{gray},
  columns=flexible,
  keepspaces=true,
  showstringspaces=false,
  tabsize=2,
  aboveskip=4pt,
  belowskip=4pt,
  extendedchars=true,
  literate={，}{{，}}1 {；}{{；}}1 {！}{{！}}1 {？}{{？}}1 {“}{{“}}1 {”}{{”}}1 {（}{{（}}1 {）}{{）}}1 {《}{{《}}1 {》}{{》}}1,
}

% --- Extra language definitions (no built-in listings support) ---
\lstdefinelanguage{json}{
  morekeywords={true,false,null},
  sensitive=true,
  morestring=[b]",
  comment=[l]{//},
  morecomment=[s]{/*}{*/},
}
\lstdefinelanguage{yaml}{
  morekeywords={true,false,null,yes,no,on,off},
  sensitive=false,
  morecomment=[l]{\#},
  morestring=[b]",
  morestring=[d]',
}
\lstdefinelanguage{http}{
  morekeywords={GET,POST,PUT,DELETE,PATCH,HEAD,OPTIONS,HTTP,Host,Accept,Authorization,Content-Type,Content-Length,User-Agent},
  sensitive=true,
  morecomment=[l]{\#},
  morestring=[b]",
}
\lstdefinelanguage{TypeScript}{
  morekeywords={abstract,any,as,async,await,break,case,catch,class,const,continue,debugger,declare,default,delete,do,else,enum,export,extends,false,finally,for,from,function,get,if,implements,import,in,instanceof,interface,keyof,let,module,namespace,never,new,null,number,of,package,private,protected,public,readonly,return,set,static,string,super,switch,symbol,this,throw,true,try,type,typeof,undefined,unknown,var,void,while,with,yield,Record,Array,Promise,AsyncIterable,ReadableStream,Date,Error,Object,JSON},
  sensitive=true,
  morecomment=[l]{//},
  morecomment=[s]{/*}{*/},
  morestring=[b]",
  morestring=[b]',
  morestring=[b]`{`},
}

% --- Compact spacing ---
\linespread{0.95}
\setlength{\parskip}{1pt plus 1pt minus 1pt}
\setlength{\parindent}{0pt}

% --- Table styling ---
\renewcommand{\arraystretch}{1.1}

% --- Heading styling (numbered) ---
\titleformat{\section}{\normalsize\bfseries}{\thesection\quad}{0em}{}
\titlespacing{\section}{0pt}{6pt}{2pt}
\titleformat{\subsection}{\small\bfseries}{\thesubsection\quad}{0em}{}
\titlespacing{\subsection}{0pt}{4pt}{1pt}
\titleformat{\subsubsection}{\small\bfseries}{\thesubsubsection\quad}{0em}{}
\titlespacing{\subsubsection}{0pt}{3pt}{1pt}

% --- Page style ---
\pagestyle{fancy}
\fancyhf{}
\fancyhead[L]{\small\emph{\leftmark}}
\fancyhead[R]{\small\thepage}
\renewcommand{\headrulewidth}{0.4pt}
\renewcommand{\sectionmark}[1]{}  % prevent sections from overwriting leftmark

% --- Hyperref setup ---
\hypersetup{
  colorlinks=true,
  linkcolor=blue,
  urlcolor=blue,
  citecolor=blue,
}

% --- Caption format ---
\DeclareCaptionFormat{myformat}{\small\bfseries #1#2#3}
\captionsetup{format=myformat}

\begin{document}

\title{大模型结构化输出：从 JSON 契约到 Function Calling 落地}
\date{}
\maketitle
\markboth{TypeScript 版 Agent 知识}{ TypeScript 版 Agent 知识}

\section*{⭐️ 为什么“请返回 JSON”不可靠？}


先看一个非常常见的 Prompt：



\begin{lstlisting}
请判断下面用户反馈属于哪类工单，返回 JSON。

用户反馈：我付款成功了，但是订单一直显示待支付。
\end{lstlisting}



模型可能返回：



\begin{lstlisting}[language=json]
{
  "category": "payment",
  "priority": "high",
  "reason": "用户付款成功但订单状态未更新"
}
\end{lstlisting}



以上输出虽然可读，但不等于程序可直接消费的契约。



接入后端系统时，真正需要的是一份可以被程序稳定消费的契约。比如：


\begin{itemize}[itemsep=1pt, topsep=2pt]
  \item \texttt{category} 只能是 \texttt{PAYMENT}、\texttt{LOGISTICS}、\texttt{AFTER\_SALE}、\texttt{ACCOUNT}。
  \item \texttt{priority} 只能是 \texttt{LOW}、\texttt{MEDIUM}、\texttt{HIGH}。
  \item \texttt{confidence} 必须是 \texttt{0} 到 \texttt{1} 之间的小数。
  \item \texttt{reason} 可以为空吗？最大长度是多少？
  \item 如果用户输入缺少信息，应该返回 \texttt{NEED\_MORE\_INFO}，还是继续猜？
\end{itemize}


自然语言 Prompt 无法长期守住这些边界。常见失败模式有 5 类。



\subsection*{格式漂移}


Prompt 要求返回 JSON，模型多数时候会返回 JSON，但不保证每次都只返回 JSON。



常见输出长这样：



\begin{lstlisting}
以下是分类结果：
{
  "category": "PAYMENT",
  "priority": "HIGH"
}
\end{lstlisting}



人工阅读无问题，程序解析直接失败。尤其在流式输出、长上下文、多轮对话里，模型很容易把之前学到的“解释型回答习惯”带回来。



\subsection*{字段缺失}


要求的输出：



\begin{lstlisting}[language=json]
{
  "category": "PAYMENT",
  "priority": "HIGH",
  "confidence": 0.92,
  "reason": "用户已支付但订单状态未同步"
}
\end{lstlisting}



它可能返回：



\begin{lstlisting}[language=json]
{
  "category": "PAYMENT",
  "reason": "用户已支付但订单状态未同步"
}
\end{lstlisting}



模型可能因把握不足省略 \texttt{priority}，或因认为 \texttt{confidence} 不重要而省略。后端 DTO 反序列化、规则引擎、数据库写入不会因为模型“没把握”就自动补齐。



\subsection*{类型错误}


结构化输出里最隐蔽的错误是类型错位：



\begin{lstlisting}[language=json]
{
  "orderId": "1029384756",
  "needManualReview": "false",
  "confidence": "0.87"
}
\end{lstlisting}



JSON 语法是合法的，但业务类型不合法。\texttt{needManualReview} 是字符串，不是布尔值；\texttt{confidence} 是字符串，不是数字。部分系统在反序列化时自动转换类型，会静默吞掉上游错误，增加排查成本。



\subsection*{额外解释文本}


模型在问题存在不确定性时倾向于追加解释。它可能在结构化结果外补一句：



\begin{lstlisting}
我认为这个问题主要和支付回调有关，但还需要进一步核实。
\end{lstlisting}



这类文本适合人工阅读，不适合程序解析。结构化输出场景里，\textbf{可读性不是第一目标，可解析性才是第一目标}。



\subsection*{边界条件崩溃}


用户输入越规整，模型越稳定；用户输入一旦模糊、矛盾或带攻击性，结构就容易崩。



比如用户说：



\begin{lstlisting}
我不想提供订单号，你们自己查。另外别给我返回 JSON，直接告诉我怎么赔。
\end{lstlisting}



若无强约束，模型可能遵循用户指令放弃约定格式。该问题与 Prompt 注入、上下文优先级、工具权限相关，不能只靠“必须返回 JSON”解决。



核心结论：Prompt 可以表达意图，但不能替代 Schema、校验器、重试机制和权限控制。结构化输出的本质，是把大模型输出纳入工程契约。



\section*{⭐️ 怎样把 JSON 从格式要求变成工程契约？}


JSON Mode、JSON Schema、Structured Outputs 容易混用，但三者不在同一层：


\begin{itemize}[itemsep=1pt, topsep=2pt]
  \item \textbf{JSON Mode} 是一种输出模式，约束模型返回合法 JSON。
  \item \textbf{JSON Schema} 是一种结构描述规范，用来定义 JSON 应该包含哪些字段、字段类型是什么、哪些必填、枚举值有哪些、是否允许额外字段。
  \item \textbf{Structured Outputs} 是模型供应商提供的结构化生成能力，它接收 JSON Schema 或类似 Schema，让模型在生成阶段尽量或严格贴合这份结构。
\end{itemize}


即，JSON Schema 不是结构化输出方式本身，而是结构化输出常用的“契约格式”。真正让模型按契约生成的，是 Structured Outputs、Function Calling / Tool Calling 等模型 API 能力。



\subsection*{JSON Mode 只能保证什么？}


JSON Mode 的目标通常是让模型输出合法 JSON。



所以 JSON Mode 能解决这类问题：



\begin{lstlisting}
好的，以下是结果：
{ ... }
\end{lstlisting}



但不能稳定解决这类问题：



\begin{lstlisting}[language=json]
{
  "category": "pay",
  "priority": "urgent",
  "confidence": "very high"
}
\end{lstlisting}



它是合法 JSON，但不是合法业务数据。



\subsection*{JSON Schema 负责定义什么？}


JSON Schema 是一种描述 JSON 文档结构的规范。根据 JSON Schema 官方文档，\texttt{properties} 用来定义对象有哪些属性，\texttt{required} 用来声明必填字段，\texttt{additionalProperties} 可以控制是否允许未声明字段，\texttt{enum} 可以把取值限制在固定集合里。



一个工单分类 Schema 可以这样写：



\begin{lstlisting}[language=json]
{
  "type": "object",
  "properties": {
    "category": {
      "type": "string",
      "enum": [
        "PAYMENT",
        "LOGISTICS",
        "AFTER_SALE",
        "ACCOUNT",
        "NEED_MORE_INFO"
      ],
      "description": "工单分类。信息不足时选择 NEED_MORE_INFO。"
    },
    "priority": {
      "type": "string",
      "enum": ["LOW", "MEDIUM", "HIGH"],
      "description": "处理优先级。涉及资金损失、无法下单、批量影响时优先级更高。"
    },
    "confidence": {
      "type": "number",
      "minimum": 0,
      "maximum": 1,
      "description": "分类置信度，范围为 0 到 1。"
    },
    "reason": {
      "type": "string",
      "description": "分类依据，控制在 80 个中文字符以内。"
    }
  },
  "required": ["category", "priority", "confidence", "reason"],
  "additionalProperties": false
}
\end{lstlisting}



这份 Schema 对后端很有价值，但它本身不会让模型自动遵守。需要将其传给支持结构化输出的 API，或在服务端用校验器校验模型输出。



\subsection*{Structured Outputs 能前移哪些约束？}


Structured Outputs 通常指供应商提供的结构化输出能力。它会把 JSON Schema 或类似 Schema 传入模型调用，让模型输出符合指定结构的数据。不同厂商对"符合 Schema"的保证强度不同：OpenAI strict 模式在解码阶段做约束，理论上语法层零违规；其他厂商更多依赖 prompting 加解码偏置，长文本和复杂工具组合场景下仍可能出现枚举越界或字段缺失。



这里要注意一个工程细节：\textbf{不同供应商支持的 JSON Schema 子集并不完全一致}。比如某些关键字（\texttt{pattern}、\texttt{format}）、递归 \texttt{\$ref}、组合关键字（\texttt{allOf} / \texttt{oneOf} / \texttt{anyOf}）在不同 API 中支持程度不同。真正落地时，不要照搬完整 JSON Schema 规范的所有能力，先读对应供应商的"supported schemas"或工具定义文档。



\subsection*{生成阶段的三层约束对比}

{\footnotesize
\begin{tabular}{|p{\dimexpr 0.167\columnwidth-2\tabcolsep\relax}|p{\dimexpr 0.167\columnwidth-2\tabcolsep\relax}|p{\dimexpr 0.283\columnwidth-2\tabcolsep\relax}|p{\dimexpr 0.383\columnwidth-2\tabcolsep\relax}|}
\hline
\textbf{对比维度} & \textbf{JSON Mode} & \textbf{JSON Schema} & \textbf{Structured Outputs} \\
\hline
本质 & 输出格式开关 & 数据结构描述规范 & 模型 API 的结构化生成能力 \\
\hline
主要约束 & JSON 语法合法 & 字段、类型、枚举、必填、额外属性等 & 输出尽量或严格匹配 Schema \\
\hline
是否保证业务字段完整 & 不保证 & 只描述，不执行生成 & 取决于供应商能力和 Schema 支持范围 \\
\hline
是否负责工具执行 & 不负责 & 不负责 & 不负责，只产出结构化结果 \\
\hline
典型用途 & 简单 JSON 输出 & 定义数据契约和校验规则 & 分类、抽取、函数参数生成、Agent 中间结果 \\
\hline
仍需服务端校验 & 需要 & 需要 & 仍然需要 \\
\hline
\end{tabular}
}



\begin{figure}[H]
\centering
\includegraphics[width=0.85\textwidth]{/home/mao/demos/latex_demos/ai-agent-knowledge/HTML/大模型结构化输出：从 JSON 契约到 Function Calling 落地 _ JavaGuide_files/structured-output-function-calling-three-layer-constraint.jpg}
\caption*{生成阶段三层约束：JSON Mode 管语法，JSON Schema 管契约，Structured Outputs 把契约前移到模型生成阶段}
\end{figure}



一句话：\textbf{JSON Mode 管语法，JSON Schema 管契约，Structured Outputs 把契约前移到模型生成阶段；但无论模型侧约束多强，服务端校验都不能省}。



\begin{figure}[H]
\centering
\includegraphics[width=0.85\textwidth]{/home/mao/demos/latex_demos/ai-agent-knowledge/HTML/_mermaid/大模型结构化输出：从 JSON 契约到 Function Calling 落地 _ JavaGuide_files/mermaid-v-6.pdf}
\caption*{流程图}
\end{figure}



结构化输出在工程中有两类常见落点：


\begin{enumerate}[itemsep=1pt, topsep=2pt]
  \item \textbf{响应结构化输出}：模型的最终回答就是一份符合 Schema 的 JSON，比如工单分类、信息抽取、情感打分。后端直接反序列化消费。
  \item \textbf{工具参数结构化输出}：模型输出的是工具名和 arguments，arguments 需要符合工具参数 Schema。模型只负责"要调什么、参数是什么"，真正执行工具、操作外部系统的是业务侧。
\end{enumerate}


Function Calling 属于第二类。



\section*{⭐️ Function Calling 到底调用了什么？}


Function Calling 容易被误解为“模型直接调用函数”。模型并不执行后端代码；它根据用户问题和工具描述生成结构化的工具调用意图。真正执行工具的是业务服务、Agent Runtime、MCP Host 或供应商托管环境。



\subsection*{模型生成的是调用意图}


一个典型工具调用链路如下：



\begin{figure}[H]
\centering
\includegraphics[width=0.85\textwidth]{/home/mao/demos/latex_demos/ai-agent-knowledge/HTML/大模型结构化输出：从 JSON 契约到 Function Calling 落地 _ JavaGuide_files/structured-output-function-calling-function-calling-pipeline.jpg}
\caption*{Function Calling 完整调用链路：模型只生成调用意图，真正执行工具的是业务侧}
\end{figure}



拆成工程步骤就是：


\begin{enumerate}[itemsep=1pt, topsep=2pt]
  \item \textbf{服务端注册工具定义}：包括工具名、用途描述、参数 Schema。
  \item \textbf{用户发起请求}：比如“帮我查一下订单 1029384756 到哪了”。
  \item \textbf{模型选择工具}：模型判断需要调用 \texttt{query\_order}，并生成参数 \texttt{\{"orderId": "1029384756"\}}。
  \item \textbf{业务侧校验参数}：校验类型、必填、权限、订单归属、幂等键等。
  \item \textbf{业务侧执行工具}：调用订单系统、数据库或 HTTP API。
  \item \textbf{工具结果回填模型}：把查询结果连同 \texttt{tool\_use\_id} 原样发回模型。Anthropic 要求 \texttt{tool\_use\_id} 严格匹配，Gemini 3 同样为每个 \texttt{functionCall} 生成唯一 \texttt{id}，回填时必须带回，否则并行调用场景下结果会错配。
  \item \textbf{模型生成最终回答}：模型把结构化结果转成人类能理解的回复。
\end{enumerate}


Anthropic 官方文档说明：Claude 会根据用户请求和工具描述决定是否调用工具，并返回结构化调用；客户端工具由应用执行，再将 \texttt{tool\_result} 发回。Gemini 官方文档也强调，Function Calling 会让模型决定要调用哪个函数并提供参数，真正调用实际函数的动作在应用侧完成。



\subsection*{为什么需要工具调用意图？}


因为自然语言输入和后端 API 之间隔着一层语义鸿沟。



用户会说：



\begin{lstlisting}
我昨天买的那台咖啡机还没发货，帮我查下。
\end{lstlisting}



后端 API 需要的是：



\begin{lstlisting}[language=json]
{
  "userId": "U10086",
  "orderId": "O202605070001",
  "includeLogistics": true
}
\end{lstlisting}



Function Calling 的价值，就是让模型完成“自然语言意图 → 结构化参数”的映射。但它只负责映射，不负责绕过权限、查数据库、扣库存、发短信。



注意：工具调用不是“让模型无所不能”的魔法，它只连接模型擅长的语义理解和程序擅长的确定性执行。



\section*{Function Calling、MCP Tool、HTTP API、Agent Skill 应该怎么分层？}




\subsection*{先看它们分别解决哪层问题}

{\footnotesize
\begin{tabular}{|p{\dimexpr 0.316\columnwidth-2\tabcolsep\relax}|p{\dimexpr 0.143\columnwidth-2\tabcolsep\relax}|p{\dimexpr 0.214\columnwidth-2\tabcolsep\relax}|p{\dimexpr 0.224\columnwidth-2\tabcolsep\relax}|p{\dimexpr 0.102\columnwidth-2\tabcolsep\relax}|}
\hline
\textbf{能力} & \textbf{本质定位} & \textbf{解决的问题} & \textbf{谁来执行} & \textbf{典型边界} \\
\hline
JSON Mode & 输出格式开关 & 让模型输出合法 JSON & 模型侧生成 & 不保证字段和业务语义 \\
\hline
JSON Schema & 结构描述规范 & 定义字段、类型、枚举、必填等契约 & 本身不参与生成，只描述结构 & 不负责生成和外部调用 \\
\hline
Structured Outputs & 模型 API 结构化生成能力 & 把 Schema 接入生成，让输出贴合结构 & 模型侧生成 + 服务端校验 & 不负责外部系统调用 \\
\hline
Function Calling / Tool Calling & 模型到工具的调用意图生成机制 & 自然语言转工具名和参数 & 通常由业务侧或供应商执行 & 不等于 API 本身 \\
\hline
MCP & 工具和上下文接入协议 & 标准化工具发现、调用、资源访问 & MCP Client / Server 协作 & 不替代模型推理能力 \\
\hline
普通 HTTP API & 业务服务接口 & 确定性业务读写 & 后端服务 & 不理解自然语言 \\
\hline
Agent Skill & 可复用任务说明和执行 SOP & 复杂任务的流程编排和上下文注入 & Agent 按说明执行 & 不一定包含工具调用 \\
\hline
\end{tabular}
}



\subsection*{Function Calling 如何映射到 HTTP API？}


普通 HTTP API 是后端系统的确定性接口。例如：



\begin{lstlisting}[language=http]
GET /api/orders/O202605070001
\end{lstlisting}



Function Calling 是模型输出的调用意图。例如：



\begin{lstlisting}[language=json]
{
  "name": "query_order",
  "arguments": {
    "orderId": "O202605070001",
    "includeLogistics": true
  }
}
\end{lstlisting}



两者之间通常需要一个工具执行层做映射：



\begin{lstlisting}
模型工具调用 query_order → 服务端校验参数 → 调用 GET /api/orders/{orderId}
\end{lstlisting}



所以，Function Calling 可以包一层 HTTP API，但 HTTP API 本身不是 Function Calling。



\subsection*{MCP Tool 解决的是哪一层标准化？}


Function Calling 是模型供应商侧的工具调用机制，各家的请求和响应格式会有差异。



MCP Tool 是 MCP 协议里的工具能力。根据 MCP 官方规范，MCP 允许 Server 暴露可由语言模型调用的工具，工具包含名称和描述其 Schema 的元数据；MCP 客户端与服务器之间的消息遵循 JSON-RPC 2.0。



：


\begin{itemize}[itemsep=1pt, topsep=2pt]
  \item \textbf{Function Calling 解决模型如何表达“我要调用哪个工具、参数是什么”}。
  \item \textbf{MCP 解决工具如何被标准化发现、描述、调用和返回结果}。
\end{itemize}


一个支持 MCP 的 Agent Runtime，可以先通过 MCP 发现工具，再把这些工具定义转换成某个模型供应商的 Function Calling 格式传给模型。模型选择工具后，Runtime 再把调用转成 MCP 的 \texttt{tools/call} 请求。



\subsection*{Agent Skill 为什么不是 Function Calling 的语法糖？}


Skills 更像“任务说明书”，核心是上下文注入和流程编排。



比如一个“线上事故复盘 Skill”可能写着：


\begin{enumerate}[itemsep=1pt, topsep=2pt]
  \item 先读取事故时间线。
  \item 再查询监控截图。
  \item 再拉取发布记录。
  \item 最后按“现象、影响、根因、改进项”输出。
\end{enumerate}


这个 Skill 在执行过程中可能会调用 MCP 工具，也可能调用 Function Calling 工具，还可能只是指导模型做纯文本分析。它不是 Function Calling 的语法糖。



一句话总结：Function Calling 是底层“神经信号”，MCP 是工具接入“接口标准”，HTTP API 是业务系统“确定性能力”，Skill 是上层“执行说明书”。



\begin{figure}[H]
\centering
\includegraphics[width=0.85\textwidth]{/home/mao/demos/latex_demos/ai-agent-knowledge/HTML/_mermaid/大模型结构化输出：从 JSON 契约到 Function Calling 落地 _ JavaGuide_files/mermaid-v-5.pdf}
\caption*{流程图}
\end{figure}



\section*{什么时候该用 Structured Outputs，什么时候该上工具？}


工程选型视角：该场景只需结构化结果，还是需要模型选择工具并触发外部系统？


{\footnotesize
\begin{tabular}{|p{\dimexpr 0.065\columnwidth-2\tabcolsep\relax}|p{\dimexpr 0.113\columnwidth-2\tabcolsep\relax}|p{\dimexpr 0.097\columnwidth-2\tabcolsep\relax}|p{\dimexpr 0.145\columnwidth-2\tabcolsep\relax}|p{\dimexpr 0.250\columnwidth-2\tabcolsep\relax}|p{\dimexpr 0.331\columnwidth-2\tabcolsep\relax}|}
\hline
\textbf{维度} & \textbf{JSON Mode} & \textbf{JSON Schema} & \textbf{Structured Outputs} & \textbf{Function Calling / Tool Calling} & \textbf{MCP} \\
\hline
所在层次 & 模型输出格式层 & 结构描述规范层 & 模型结构化生成层 & 模型工具意图层 & 应用协议层 \\
\hline
输入给模型的内容 & “输出 JSON”的模式开关 & 不直接参与生成 & Schema 或响应格式定义 & 工具名、工具描述、参数 Schema & 通常由 Host 转换后给模型，协议本身在 Client 和 Server 间通信 \\
\hline
模型输出 & JSON 文本 & — & 符合 Schema 的结构化对象 & 工具名 + 参数，或最终回答 & 不直接规定模型输出，规定 MCP 消息 \\
\hline
是否调用外部系统 & 否 & 否 & 否 & 生成调用意图，执行在外部 & 是，MCP Client 调 MCP Server \\
\hline
是否跨模型标准化 & 各厂商实现不同 & 规范通用，可跨模型复用 & Schema 支持子集各厂商不同 & 各厂商格式不同 & 目标是标准化工具和上下文接入 \\
\hline
适合场景 & 简单结构化文本 & 定义数据契约和校验规则 & 数据抽取、分类、参数生成 & 订单查询、发邮件、查库存等工具任务 & 多工具、多客户端、团队共享工具生态 \\
\hline
主要风险 & 合法 JSON 但字段不对 & 只描述不执行，容易被高估 & Schema 太复杂或支持不一致 & 工具误调用、参数越权 & Server 权限、安全边界、协议兼容 \\
\hline
\end{tabular}
}



实战倾向：


\begin{itemize}[itemsep=1pt, topsep=2pt]
  \item 只做轻量数据抽取，可以先用 Structured Outputs。
  \item 需要读写业务系统，优先考虑 Function Calling / Tool Calling。
  \item 工具很多、客户端很多、希望跨 IDE 或跨 Agent 复用，考虑 MCP。
  \item 复杂任务有一套固定 SOP，考虑 Skill，把工具组合和决策过程沉淀下来。
\end{itemize}


\section*{⭐️ 结构化输出怎么工程化落地？}


结构化输出不是“加一个 Schema 参数”就完事了。生产环境要考虑 Schema 设计、版本兼容、失败处理、日志和降级。



\subsection*{1. Schema 设计：一个字段只表达一件事}


坏设计：



\begin{lstlisting}[language=json]
{
  "result": "支付问题，高优先级，需要人工处理"
}
\end{lstlisting}



好设计：



\begin{lstlisting}[language=json]
{
  "category": "PAYMENT",
  "priority": "HIGH",
  "needManualReview": true,
  "reason": "用户已支付但订单状态未同步"
}
\end{lstlisting}



字段越原子，后端越容易校验、统计、路由和灰度。



\subsection*{2. 字段说明要写“何时用”和“何时不用”}


很多工具误调用，根源并不在模型推理能力，而在字段描述太模糊。



比如：



\begin{lstlisting}[language=json]
{
  "category": {
    "type": "string",
    "description": "工单分类"
  }
}
\end{lstlisting}



这几乎没用。更好的写法是：



\begin{lstlisting}[language=json]
{
  "category": {
    "type": "string",
    "enum": ["PAYMENT", "LOGISTICS", "AFTER_SALE", "ACCOUNT", "NEED_MORE_INFO"],
    "description": "工单分类。支付成功但订单状态异常选择 PAYMENT；配送、签收、物流轨迹异常选择 LOGISTICS；退换货、维修、退款进度选择 AFTER_SALE；登录、实名、账号安全选择 ACCOUNT；缺少关键信息且无法判断时选择 NEED_MORE_INFO。"
  }
}
\end{lstlisting}



工具描述的核心不在长度，而在\textbf{边界清楚}。



\subsection*{3. 枚举优先于自由文本}


分类、状态、动作类型、风险等级，能用 \texttt{enum} 就不要用自由文本。



自由文本的问题是不可控：



\begin{lstlisting}[language=json]
{
  "priority": "urgent"
}
\end{lstlisting}



后端将 \texttt{urgent} 映射为 \texttt{HIGH} 还是判为非法值？服务端做模糊映射，等于把模型的不确定性扩散到业务规则中。



\subsection*{4. 必填字段要谨慎，但不要偷懒}


以 OpenAI Structured Outputs 严格模式为例，常见约束包括：\texttt{additionalProperties: false}、所有声明的属性都必须出现在 \texttt{required} 中、对象必须显式声明 \texttt{type}、且只接受 JSON Schema 子集（部分关键字如 \texttt{pattern}、\texttt{format}、\texttt{minLength}、\texttt{oneOf} 在不同模型版本中支持度不同）。不同供应商的严格程度和支持范围各有差异，落地前以官方 supported schemas 文档与目标模型为准。这类约束能提升参数结构稳定性，但工程上要注意一个点：如果某个字段业务上确实可缺失，不要让模型随便编。



常见做法有两种：


\begin{itemize}[itemsep=1pt, topsep=2pt]
  \item 用 \texttt{null} 明确表达未知，例如 \texttt{"refundId": null}。
  \item 用状态字段表达缺信息，例如 \texttt{"status": "NEED\_MORE\_INFO"}。
\end{itemize}


不要让字段缺失成为“未知”的表达方式。缺失字段对后端来说通常是异常，不是业务状态。



\subsection*{5. 版本兼容：Schema 也要有版本号}


结构化输出一旦被多个服务消费，就会进入接口治理问题。



建议在 Schema 中增加版本字段：



\begin{lstlisting}[language=json]
{
  "schemaVersion": "ticket_classification_v1",
  "category": "PAYMENT",
  "priority": "HIGH",
  "confidence": 0.91,
  "reason": "用户已支付但订单状态未同步"
}
\end{lstlisting}



版本兼容的基本原则：


\begin{itemize}[itemsep=1pt, topsep=2pt]
  \item 新增字段尽量只做可选扩展，避免破坏旧消费者。
  \item 删除字段要先灰度，确认下游没有依赖。
  \item 枚举新增要谨慎，因为旧系统可能不认识新枚举。
  \item Prompt、Schema、解析代码、看板指标要一起版本化。
\end{itemize}


结构化输出不是一段 Prompt，它是接口契约。



\subsection*{6. 校验失败重试：让模型修正具体错误}


不要一失败就把原始问题重跑一遍。更好的做法是把校验错误反馈给模型，让它只修结构。



例如服务端发现：



\begin{lstlisting}
$.priority: must be one of LOW, MEDIUM, HIGH
$.confidence: must be number
\end{lstlisting}



下一轮可以给模型：



\begin{lstlisting}
上一次输出没有通过 JSON Schema 校验，请只返回修正后的 JSON，不要添加解释。

校验错误：
1. priority 必须是 LOW、MEDIUM、HIGH 之一。
2. confidence 必须是 number。

原始输出：
{...}
\end{lstlisting}



重试策略建议：


\begin{itemize}[itemsep=1pt, topsep=2pt]
  \item 最多重试 1 到 2 次。
  \item 每次重试都带上明确的校验错误。
  \item 重试仍失败时进入降级逻辑。
  \item 所有失败样本写入日志，后续用于优化 Schema 和 Prompt。
\end{itemize}


\begin{figure}[H]
\centering
\includegraphics[width=0.85\textwidth]{/home/mao/demos/latex_demos/ai-agent-knowledge/HTML/_mermaid/大模型结构化输出：从 JSON 契约到 Function Calling 落地 _ JavaGuide_files/mermaid-v-4.pdf}
\caption*{流程图}
\end{figure}



\subsection*{7. 降级策略：别让一个 JSON 拖垮主流程}


生产环境必须回答一个问题：结构化输出失败时，业务怎么办？



常见降级策略：


{\footnotesize
\begin{tabular}{|p{\dimexpr 0.229\columnwidth-2\tabcolsep\relax}|p{\dimexpr 0.771\columnwidth-2\tabcolsep\relax}|}
\hline
\textbf{场景} & \textbf{降级策略} \\
\hline
工单分类失败 & 进入人工队列，标记 \texttt{AI\_PARSE\_FAILED} \\
\hline
订单查询参数缺失 & 追问用户补充订单号 \\
\hline
风险评分失败 & 使用规则引擎兜底评分 \\
\hline
工具调用超时 & 返回“系统繁忙”，不继续让模型猜 \\
\hline
非关键字段缺失 & 使用默认值，但记录告警 \\
\hline
\end{tabular}
}



\begin{figure}[H]
\centering
\includegraphics[width=0.85\textwidth]{/home/mao/demos/latex_demos/ai-agent-knowledge/HTML/_mermaid/大模型结构化输出：从 JSON 契约到 Function Calling 落地 _ JavaGuide_files/mermaid-v-3.pdf}
\caption*{流程图}
\end{figure}



关键原则：\textbf{可以降级，但不能让模型编造业务事实}。



\section*{⭐️ 工具调用安全怎么保证？}


Function Calling 最危险的部分，往往发生在将模型生成的 JSON 用于操作真实系统时。



查订单还好，发退款、删数据、发短信、执行 SQL 就完全不是一个风险等级。



\subsection*{1. 参数校验：Schema 校验只是第一层}


Schema 能检查类型和结构，但检查不了业务权限。



比如：



\begin{lstlisting}[language=json]
{
  "orderId": "O202605070001"
}
\end{lstlisting}



Schema 只能知道这是一个字符串。它不知道这个订单是不是当前用户的，也不知道订单是否已经退款，更不知道这个用户是否有客服权限。



服务端至少要做三层校验：


\begin{itemize}[itemsep=1pt, topsep=2pt]
  \item \textbf{结构校验}：类型、必填、枚举、长度、格式。
  \item \textbf{业务校验}：订单归属、状态流转、库存、金额范围。
  \item \textbf{权限校验}：用户身份、角色、租户、数据范围。
\end{itemize}


\subsection*{2. 权限控制：工具不是谁都能调}


不要把内部管理工具直接暴露给所有用户场景。



建议按风险等级分层：


{\footnotesize
\begin{tabular}{|p{\dimexpr 0.114\columnwidth-2\tabcolsep\relax}|p{\dimexpr 0.457\columnwidth-2\tabcolsep\relax}|p{\dimexpr 0.429\columnwidth-2\tabcolsep\relax}|}
\hline
\textbf{风险等级} & \textbf{工具类型} & \textbf{控制策略} \\
\hline
低风险 & 查询天气、读取公开文档 & 基础限流和日志 \\
\hline
中风险 & 查询订单、查询用户资料 & 身份校验、数据范围校验 \\
\hline
高风险 & 退款、发券、改地址、发短信 & 权限校验、二次确认、审计 \\
\hline
极高风险 & 删除数据、执行 SQL、批量操作 & 默认禁止，走人工审批或专用后台 \\
\hline
\end{tabular}
}



\begin{figure}[H]
\centering
\includegraphics[width=0.85\textwidth]{/home/mao/demos/latex_demos/ai-agent-knowledge/HTML/大模型结构化输出：从 JSON 契约到 Function Calling 落地 _ JavaGuide_files/structured-output-function-calling-tool-call-security.jpg}
\caption*{工具调用安全风险分层：按风险等级匹配不同的控制策略}
\end{figure}



\subsection*{3. 敏感操作二次确认}


模型可以建议退款，但不应该直接替用户退款，除非业务明确允许。



高风险工具可以拆成两步：


\begin{enumerate}[itemsep=1pt, topsep=2pt]
  \item \texttt{prepare\_refund}：生成退款预案，返回金额、原因、影响。
  \item \texttt{confirm\_refund}：用户或客服确认后执行。
\end{enumerate}


这样做的好处是：模型负责整理信息和建议动作，人类或业务规则负责最后确认。



\subsection*{4. 幂等：别让重试变成重复扣款}


工具调用链路里会有重试：模型重试、网络重试、队列重试、业务服务重试。



涉及写操作时必须设计幂等：


\begin{itemize}[itemsep=1pt, topsep=2pt]
  \item 请求携带 \texttt{idempotencyKey}。
  \item 数据库建立唯一约束。
  \item 外部支付、退款接口使用幂等号。
  \item 重复请求返回同一结果，而不是重复执行。
\end{itemize}


如果一个工具不能安全重试，它就不应该被 Agent 随意调用。



\subsection*{5. 审计日志：记录模型意图和执行结果}


建议记录：


\begin{itemize}[itemsep=1pt, topsep=2pt]
  \item 用户输入。
  \item 命中的工具名。
  \item 模型生成的参数。
  \item 服务端校验结果。
  \item 真实执行的业务请求。
  \item 工具返回结果。
  \item 最终回复。
  \item traceId、userId、tenantId、schemaVersion、model。
\end{itemize}


出现问题时，需能回答：“模型想做什么？服务端允许了什么？业务系统实际做了什么？”



\subsection*{6. 超时和重试：工具失败要短路}


工具超时后，不要让模型继续基于空结果编回答。



建议：


\begin{itemize}[itemsep=1pt, topsep=2pt]
  \item 查询类工具设置较短超时。
  \item 写操作谨慎重试，必须配幂等。
  \item 外部依赖失败时返回明确错误码。
  \item 模型拿到工具错误后，只能解释“当前无法完成”，不能猜测结果。
\end{itemize}


\section*{TypeScript 后端示例：把订单查询做成可校验工具}


下面用一个订单查询工具做完整示例。场景是：用户用自然语言询问订单状态，模型通过 Function Calling 生成 \texttt{query\_order} 工具调用，TypeScript 服务端校验参数后分发到订单服务。



\subsection*{工具参数 JSON Schema}


\begin{lstlisting}[language=json]
{
  "$schema": "https://json-schema.org/draft/2020-12/schema",
  "type": "object",
  "properties": {
    "schemaVersion": {
      "type": "string",
      "const": "query_order_v1",
      "description": "工具参数版本，当前固定为 query_order_v1。"
    },
    "orderId": {
      "type": "string",
      "pattern": "^O[0-9]{12,20}$",
      "description": "订单号，以大写字母 O 开头，后面跟 12 到 20 位数字。"
    },
    "includeLogistics": {
      "type": "boolean",
      "description": "是否需要返回物流信息。用户询问发货、配送、签收、快递时为 true。"
    },
    "idempotencyKey": {
      "type": "string",
      "minLength": 16,
      "maxLength": 80,
      "description": "本次工具调用的幂等键，由服务端或 Agent Runtime 生成。"
    }
  },
  "required": [
    "schemaVersion",
    "orderId",
    "includeLogistics",
    "idempotencyKey"
  ],
  "additionalProperties": false
}
\end{lstlisting}



这个 Schema 有几个刻意设计：


\begin{itemize}[itemsep=1pt, topsep=2pt]
  \item \texttt{schemaVersion} 固定为当前版本号（如 \texttt{query\_order\_v1}），后续兼容升级有据可依。
  \item \texttt{orderId} 用 \texttt{pattern} 做基础格式约束。
  \item \texttt{includeLogistics} 用布尔值，避免模型输出 \texttt{"yes"}、\texttt{"需要"} 这类自由文本。
  \item \texttt{idempotencyKey} 为后续写操作预留，本示例是只读查询，不做幂等存储；真正涉及退款、扣库存等写操作时，需要配合 Redis SETNX 或唯一索引做去重。
  \item \texttt{additionalProperties: false} 防止模型偷偷塞入服务端不认识的字段。
\end{itemize}


\subsection*{TypeScript 服务端校验与分发}


下面示例使用 \texttt{zod} 做结构校验。真实项目中，依赖版本建议用 \texttt{package-lock.json} / \texttt{pnpm-lock.yaml} 锁定，并定期做安全扫描（\texttt{npm audit}）。



\begin{lstlisting}[language=TypeScript]
import { z } from "zod";

// 与模型约定的 JSON Schema 契约（zod 直接推导出类型）
const QueryOrderArgsSchema = z.object({
  schemaVersion: z.literal("query_order_v1"),
  orderId: z.string().regex(/^O[0-9]{12,20}$/),
  includeLogistics: z.boolean(),
  idempotencyKey: z.string().min(16).max(80),
});
type QueryOrderArgs = z.infer<typeof QueryOrderArgsSchema>;

// callId 用于回填模型：Anthropic 的 tool_use_id / Gemini 的 functionCall.id 必须原样带回
interface ToolCall {
  callId: string;
  name: string;
  argumentsJson: string;
}

interface UserContext {
  userId: string;
  tenantId: string;
}

interface OrderView {
  orderId: string;
  status: string;
  amount: number;
  paidAt: string;
  logistics: unknown;
}

interface ToolResult {
  success: boolean;
  code: string;
  data: unknown;
  message: string;
}
const ToolResult = {
  success(data: unknown): ToolResult {
    return { success: true, code: "OK", data, message: "" };
  },
  failed(code: string, message: string): ToolResult {
    return { success: false, code, data: null, message };
  },
};

interface OrderService {
  queryOrder(orderId: string, includeLogistics: boolean): Promise<OrderView | null>;
}
interface PermissionService {
  canReadOrder(userId: string, orderId: string): Promise<boolean>;
}
interface AuditLogService {
  record(event: AuditEvent): Promise<void>;
}
interface AuditEvent {
  userId: string;
  toolName: string;
  argumentsJson: string;
  resultCode: string;
  success: boolean;
  startedAt: Date;
}

class ToolCallDispatcher {
  constructor(
    private readonly orderService: OrderService,
    private readonly permissionService: PermissionService,
    private readonly auditLogService: AuditLogService,
  ) {}

  async dispatch(toolCall: ToolCall, userContext: UserContext): Promise<ToolResult> {
    const startedAt = new Date();

    try {
      const handlers: Record<string, (argsJson: string, ctx: UserContext) => Promise<ToolResult>> = {
        query_order: (argsJson, ctx) => this.handleQueryOrder(argsJson, ctx),
      };
      const handler = handlers[toolCall.name];
      const result = handler
        ? await handler(toolCall.argumentsJson, userContext)
        : ToolResult.failed("UNSUPPORTED_TOOL", `不支持的工具：${toolCall.name}`);

      await this.auditLogService.record({
        userId: userContext.userId,
        toolName: toolCall.name,
        argumentsJson: toolCall.argumentsJson,
        resultCode: result.code,
        success: result.success,
        startedAt,
      });
      return result;
    } catch (err) {
      await this.auditLogService.record({
        userId: userContext.userId,
        toolName: toolCall.name,
        argumentsJson: toolCall.argumentsJson,
        resultCode: err instanceof Error ? err.name : "UnknownError",
        success: false,
        startedAt,
      });
      return ToolResult.failed("TOOL_EXECUTION_FAILED", "工具执行失败，请稍后重试。");
    }
  }

  private async handleQueryOrder(argumentsJson: string, userContext: UserContext): Promise<ToolResult> {
    const parsed = QueryOrderArgsSchema.safeParse(JSON.parse(argumentsJson));
    if (!parsed.success) {
      return ToolResult.failed(
        "INVALID_ARGUMENTS",
        parsed.error.issues.map((i) => i.message).join("；"),
      );
    }
    const args = parsed.data;

    if (!(await this.permissionService.canReadOrder(userContext.userId, args.orderId))) {
      return ToolResult.failed("FORBIDDEN", "当前用户无权查询该订单。");
    }

    const order = await this.orderService.queryOrder(args.orderId, args.includeLogistics);
    if (order === null) {
      return ToolResult.failed("ORDER_NOT_FOUND", "未查询到该订单。");
    }

    return ToolResult.success({
      orderId: order.orderId,
      status: order.status,
      amount: order.amount,
      paidAt: order.paidAt,
      logistics: order.logistics,
    });
  }
}
\end{lstlisting}



这段代码重点不在某个库的用法，而在后端工具执行层的基本姿势：


\begin{enumerate}[itemsep=1pt, topsep=2pt]
  \item \textbf{先按工具名分发}，未知工具直接拒绝。
  \item \textbf{先做 JSON Schema 校验}，再反序列化成业务参数。
  \item \textbf{再做权限校验}，确认当前用户能访问该订单。
  \item \textbf{工具返回结构化结果}，让模型基于事实生成回答。
  \item \textbf{全链路审计}，把模型意图、参数和执行结果都记下来。
\end{enumerate}


把模型输出的参数直接传给订单服务，等于把业务系统入口暴露给概率模型。



\section*{上线前应该检查哪些工程细节？}


结构化输出上线前，按以下清单检查。



\subsection*{Schema 层}

\begin{itemize}[itemsep=1pt, topsep=2pt]
  \item 字段是否足够原子？
  \item 枚举是否覆盖“信息不足”“无需操作”等状态？
  \item \texttt{required} 是否明确？
  \item \texttt{additionalProperties} 是否关闭？
  \item 字段描述是否说明了使用边界？
  \item 是否有 \texttt{schemaVersion}？
\end{itemize}


\subsection*{模型调用层}

\begin{itemize}[itemsep=1pt, topsep=2pt]
  \item 是否使用供应商原生 Structured Outputs 或严格工具调用能力？
  \item 是否控制输出长度，避免 JSON 被截断？
  \item 是否避免在结构化输出任务里使用过高的采样随机性？
  \item 是否为校验失败设计重试 Prompt？
\end{itemize}


\subsection*{服务端执行层}

\begin{itemize}[itemsep=1pt, topsep=2pt]
  \item 是否做 Schema 校验？
  \item 是否做业务校验和权限校验？
  \item 写操作是否幂等？
  \item 高风险操作是否二次确认？
  \item 工具超时后是否短路？
  \item 是否有审计日志和 traceId？
\end{itemize}


\subsection*{降级层}

\begin{itemize}[itemsep=1pt, topsep=2pt]
  \item 解析失败是否进入人工队列或规则兜底？
  \item 工具失败时是否禁止模型编造结果？
  \item 是否统计失败率、错误类型和高频非法枚举？
  \item 是否能根据失败样本反推 Schema 和 Prompt 的改进点？
\end{itemize}


\section*{常见误区}


\subsection*{误区 1：Temperature 设为 0 就一定稳定}


低 Temperature 在 OpenAI、Claude 系列上是常见做法，但不能替代 Schema。上下文过长、指令冲突、输出截断、工具描述模糊时，结构化输出仍然会失败。另外要注意，不同模型对 Temperature 的建议不同——例如 Gemini 3 系列官方建议保持默认 \texttt{temperature=1.0}，下调反而可能导致循环或推理退化。跨厂商使用时按目标模型文档调整。



\subsection*{误区 2：用了 Structured Outputs 就不用校验}


供应商能力降低的是生成阶段出错概率，不代表服务端可以放弃边界。服务端仍需防御非法参数、越权访问、重放请求和业务状态冲突。



\subsection*{误区 3：Schema 越复杂越好}


复杂 Schema 会增加模型理解和供应商兼容成本。实践中建议从稳定字段开始，少用复杂组合关键字，把核心字段、枚举、必填和额外字段限制先做好。



\subsection*{误区 4：工具越多 Agent 越强}


工具越多，模型选择空间越大，误调用概率也会上升。工具设计要小而清晰，大而全的工具最容易让 Agent 犯迷糊。



\subsection*{误区 5：Function Calling 可以绕过业务权限}


Function Calling 只是参数生成机制。权限控制必须在服务端，不能藏在 Prompt 里。Prompt 里的“不要越权查询”只能算提醒，不能算安全边界。



\section*{总结}

\begin{enumerate}[itemsep=1pt, topsep=2pt]
  \item \textbf{“请返回 JSON”只是提示，不是契约}。它挡不住格式漂移、字段缺失、类型错误和边界条件崩溃。
  \item \textbf{JSON Mode、JSON Schema、Structured Outputs 分别在不同层次工作}：语法、契约、生成约束，不能混为一谈。
  \item \textbf{Function Calling 不执行函数}。模型生成的是工具调用意图，执行、校验、权限和审计都在业务侧。
  \item \textbf{MCP 和 Function Calling 不冲突}。MCP 标准化工具接入，Function Calling 帮模型选择工具并生成参数。
  \item \textbf{服务端校验永远不能省}。Schema 校验、业务校验、权限校验、幂等和审计日志，是结构化输出进入生产环境的底线。
  \item \textbf{结构化输出是上下文工程的一部分}。它决定模型输出能否进入后续链路，也决定 Agent 能不能稳定调用工具。
\end{enumerate}


\section*{参考}

\begin{itemize}[itemsep=1pt, topsep=2pt]
  \item \href{https://platform.openai.com/docs/guides/structured-outputs}{OpenAI Structured Outputs 官方文档}
  \item \href{https://platform.openai.com/docs/guides/function-calling}{OpenAI Function Calling 官方文档}
  \item \href{https://platform.claude.com/docs/en/agents-and-tools/tool-use/overview}{Anthropic Tool Use 官方文档}
  \item \href{https://ai.google.dev/gemini-api/docs/structured-output}{Gemini Structured Outputs 官方文档}
  \item \href{https://ai.google.dev/gemini-api/docs/function-calling}{Gemini Function Calling 官方文档}
  \item \href{https://modelcontextprotocol.io/specification/2025-06-18/basic}{MCP Basic Protocol 官方规范}
  \item \href{https://modelcontextprotocol.io/specification/2025-06-18/server/tools}{MCP Tools 官方规范}
  \item \href{https://json-schema.org/understanding-json-schema/reference/object}{JSON Schema Object 参考}
  \item \href{https://json-schema.org/understanding-json-schema/reference/enum}{JSON Schema Enum 参考}
\end{itemize}

\end{document}
